\section{本章小结}
\label{sec:cone-summary}

章首的三件事现在有了各自的归属。带预算的最小二乘是凸二次规划，距离上限的交集是二阶锥规划，相关系数的可行范围由半正定约束决定。把问题写成这些标准形式，是为了把目标的曲率、可行集的形状和可检验的对偶界摆到明面上。

二次规划要看的是$\bm{Q}$在可行差方向上的曲率，而不是全空间。这一曲率决定凸性和唯一性；凸时KKT条件是全局证书，非凸时KKT点可以是局部极小点或鞍点，证书不再成立。等式约束的情形化为线性方程，方程可逆只保证驻点唯一，极小性还要看约化Hessian。

锥规划把非负约束换成闭凸锥，乘子换成对偶锥的元素。弱对偶无条件成立，给出的下界可以直接验证；强对偶和最优解的达到则需要严格可行性等条件，线性规划的结论不能原样搬到半正定规划。

表~\ref{tab:cone-summary}比较后半章的两种操作。

\begin{BookInlineTable}
  \centering
  \begin{tabular}{lll}
    \toprule
    & 等价改写 & 半正定松弛\\
    \midrule
    可行集 & 投影后与原集合相同 & 严格扩大（一般情形）\\
    最优值 & 相等 & 不超过原最优值\\
    恢复原变量 & 丢弃辅助变量 & 秩一时直接读出，否则取整\\
    所需核对 & 改写的两个方向 & 可行性、对偶可行性、间隙\\
    典型工具 & 上图形式、Schur补 & 提升、删秩、随机取整\\
    \bottomrule
  \end{tabular}
  \caption{等价改写与半正定松弛的对比。}
  \label{tab:cone-summary}
\end{BookInlineTable}

算法返回的矩阵、原问题的可行解和最优性证书是三样不同的东西，需要分别检查。松弛最优矩阵不一定对应可行点，近似的原始值不是下界，只有经原约束核对的点给出上界，只有对偶可行的解给出下界，两者之差才是可以报告的结论。
